\section{Summary}

In this paper, a simple new quark-field smearing has been proposed and tested in
numerical investigations using QCD gauge fields with $N_f=2+1$ dynamical
flavors. There is substantial freedom in the construction of creation operators
for hadrons.  This has been exploited in this work to define an algorithm
that applies a low-rank projection operator to the quark fields of the path
integral on each time-slice in order to define smooth modes for subsequent
operator construction.  This makes affordable the evaluation of all
information about quark propagation required to measure distilled hadronic
correlation functions.  Once the propagation matrix has been evaluated,
correlations between arbitrary choices of creation and annihilation
operators can be determined without requiring any further inversions. The
theoretical framework needed to extend the scope of measurements using this
technology to include matrix elements was described briefly.

Some initial tests of the method demonstrate that restricting quark fields to
lie in a very low-dimensional space of smooth fields does not substantially
degrade the quality of Monte Carlo evaluations of correlation functions and
in
many cases, statistical accuracy is enhanced. The most substantial drawback,
still to be overcome, is the growth in cost of the algorithm as the lattice
volume in physical units increases. Since the method is concerned with
long-distance physics, no extra difficulties in approaching the continuum
limit are anticipated, but this would need to be confirmed by practical
testing.

The method is currently in its infancy and a number of research directions are
outlined in this paper. The primary focus of the development efforts of this
collaboration is to include stochastic evaluation of correlation
functions in an optimal way and results will appear soon. The collaboration
will use the techniques outlined in this paper in a range of spectroscopy
determinations in the near future.
